\chapter{Design Development}
\label{ch:design}

\begin{guidance}[What this chapter does]
Generates candidate solutions, evaluates them against criteria derived from
the specification, and justifies the chosen solution. The Handbook also asks
you to consider the environmental and societal impact of the solution
(learning outcome R3). Generate several concepts -- a single concept is
usually a weak point. Typical length: 1,200--1,800 words.

Methods for generating and evaluating concepts: \textcite{cross2008},
\textcite{dym2009}; concept selection matrices: \textcite{ulrich2012}.
\end{guidance}

\section{Candidate solutions}
\label{sec:concepts}

\begin{guidance}
Sketches or block diagrams with a short description of each concept and the
enabling technologies. Keep them at the same level of detail so that the
comparison is fair.
\end{guidance}

\begin{figure}[htbp]
  \centering
  \begin{subfigure}[b]{0.32\linewidth}
    \placeholder{Concept A sketch}
    \caption{Concept A}
  \end{subfigure}\hfill
  \begin{subfigure}[b]{0.32\linewidth}
    \placeholder{Concept B sketch}
    \caption{Concept B}
  \end{subfigure}\hfill
  \begin{subfigure}[b]{0.32\linewidth}
    \placeholder{Concept C sketch}
    \caption{Concept C}
  \end{subfigure}
  \caption{Candidate solutions (replace the placeholders with your sketches)}
  \label{fig:concepts}
\end{figure}

\section{Options analysis}
\label{sec:options}

\begin{guidance}
Derive the criteria from the specification, justify each weight, score each
concept (e.g.\ 1--5) and compare the weighted totals. The macros below do the
arithmetic: \verb|\resetscores{A,B,C}| at the start, \verb|\weight{...}| for
each criterion, \verb|\score{A}{...}| for each concept and
\verb|\weightsum| and \verb|\scoretotal{A}| in the last row. Use a plain
\texttt{tabular} here, not \texttt{tabularx}. Discuss the result: is the
winner clear, and how sensitive is it to the weights?
\end{guidance}

\begin{table}[htbp]
  \centering
  \caption{Weighted decision matrix (scores 1--5, example)}
  \label{tab:decision}
  \resetscores{A,B,C}
  \begin{tabular}{llccc}
    \toprule
    Criterion (requirement) & Weight & Concept A & Concept B & Concept C \\
    \midrule
    Accuracy (\ref{req:accuracy})    & \weight{0.30} & \score{A}{4} & \score{B}{5} & \score{C}{3} \\
    Battery life (\ref{req:battery}) & \weight{0.25} & \score{A}{3} & \score{B}{4} & \score{C}{5} \\
    Cost (\ref{req:cost})            & \weight{0.25} & \score{A}{5} & \score{B}{2} & \score{C}{4} \\
    Size (\ref{req:size})            & \weight{0.10} & \score{A}{3} & \score{B}{3} & \score{C}{4} \\
    Ease of manufacture               & \weight{0.10} & \score{A}{4} & \score{B}{3} & \score{C}{3} \\
    \midrule
    \textbf{Weighted total} & \weightsum & \textbf{\scoretotal{A}} & \textbf{\scoretotal{B}} & \textbf{\scoretotal{C}} \\
    \bottomrule
  \end{tabular}
\end{table}

\section{Selected solution}
\label{sec:selected}

\begin{guidance}
Justify the choice with reference to the matrix, the specification and the
objectives. Then describe the chosen design in enough detail for the
implementation chapter: system architecture, key calculations, component
selection.
\end{guidance}

\section{Environmental and societal impact}
\label{sec:impact}

\begin{guidance}
Materials, energy use, end of life and recyclability, safety of users,
accessibility, sustainability; how the design minimises adverse effects.
\end{guidance}
